% Figure: reheat Rankine cycle on the T-s plane. After partial expansion in the
% high-pressure turbine, the steam returns to the boiler for reheat back to a
% high temperature, then completes its expansion in the low-pressure turbine.
% The two-step expansion keeps the turbine-exit quality higher than a single
% expansion to the same condenser pressure would.
\begin{figure}[htbp]
\centering
\begin{tikzpicture}
\begin{axis}[
  width=12cm, height=8.5cm,
  xlabel={specific entropy $s$ (kJ/kg$\cdot$K)},
  ylabel={temperature $T$ ($^{\circ}$C)},
  xmin=0, xmax=8.5, ymin=0, ymax=650,
  axis lines=left,
  tick label style={font=\scriptsize}, label style={font=\small},
  clip=false,
  ytick={0,100,200,300,400,500,600},
]
% === saturation dome (schematic): liquid line up the left, vapor line down
% the right, meeting at the critical point near (4.4, 374). ===
\addplot[enggray, thick, smooth] coordinates {
  (0.0,5) (1.3,100) (2.2,200) (3.0,300) (3.8,360) (4.4,374)
};
\addplot[enggray, thick, smooth] coordinates {
  (4.4,374) (5.5,360) (6.8,300) (7.6,200) (8.1,100) (8.3,46)
};
\node[font=\scriptsize, enggray, anchor=north] at (axis cs:4.4,374)
  {critical point};
% === the cycle states ===
% 1 pump inlet (sat liquid at condenser, ~46 C), 2 pump exit (~47 C),
% 3 boiler exit / HP turbine inlet (superheated, ~500 C),
% 4 HP turbine exit (reheat inlet, ~300 C, partly expanded),
% 5 reheat exit / LP turbine inlet (~500 C, higher s),
% 6 LP turbine exit (condenser pressure, two-phase ~46 C, high quality).
\addplot[engblue, very thick] coordinates {
  (0.6,46) (0.65,48)
};
\addplot[engred, very thick] coordinates {
  (0.65,48) (2.0,260) (3.0,360) (6.0,500)
};
\addplot[engblue, very thick] coordinates {
  (6.0,500) (6.3,300)
};
\addplot[engred, very thick] coordinates {
  (6.3,300) (7.5,500)
};
\addplot[engblue, very thick] coordinates {
  (7.5,500) (7.9,46)
};
\addplot[engblue, very thick] coordinates {
  (7.9,46) (0.6,46)
};
% labels
\node[font=\scriptsize, anchor=east] at (axis cs:0.6,46) {1,2};
\node[font=\scriptsize, anchor=south] at (axis cs:6.0,500) {3};
\node[font=\scriptsize, anchor=north east] at (axis cs:6.3,300) {4};
\node[font=\scriptsize, anchor=south] at (axis cs:7.5,500) {5};
\node[font=\scriptsize, anchor=west] at (axis cs:7.9,46) {6};
\node[font=\small, engred, anchor=west] at (axis cs:3.1,430)
  {boiler + reheat};
\node[font=\small, engblue, anchor=north west] at (axis cs:6.6,420)
  {HP / LP turbine};
% mark the high-quality exit advantage
\draw[-{Latex[length=2mm]}, enggray, thick]
  (axis cs:8.3,150) -- (axis cs:7.95,60);
\node[font=\scriptsize, enggray, anchor=west] at (axis cs:7.5,170)
  {exit kept dry};
\end{axis}
\end{tikzpicture}
\caption[Reheat Rankine cycle on the temperature-entropy plane]{The reheat
Rankine cycle on the $T$-$s$ plane, drawn schematically against the saturation
dome (grey). Water is pumped (1 to 2), heated and superheated in the boiler
(2 to 3), and expanded in the high-pressure turbine to an intermediate
pressure (3 to 4). It then returns to the boiler for reheat back to a high
temperature (4 to 5) before completing its expansion in the low-pressure
turbine (5 to 6) and condensing (6 to 1). Splitting the expansion and
reheating between the halves keeps the turbine-exit state (6) at higher
quality than a single expansion from state 3 to the condenser pressure would
reach, which limits the droplet erosion that wet steam inflicts on the last
turbine stages. Reheat also raises the average temperature of heat addition,
so it improves efficiency as well as protecting the blades.}
\label{fig:vol04ch03:reheat-rankine-ts}
\end{figure}
